\chapter{Chốt Tiết, Chốt Tuần}
\chapterintro{Câu hỏi cuối giờ để lại dư vị}{Một câu hỏi chốt đúng giúp bài học không rơi xuống khi tiếng trống vang lên.}

\section{Hôm nay em mang về một câu hay một ý?}
\begin{cauhoibox}{Câu hỏi}
Hôm nay em mang về một câu hay một ý?
\end{cauhoibox}
\begin{taisaocoich}
Câu hỏi này đủ mở để học sinh trả lời theo cách riêng, nhưng vẫn đủ gọn để chốt nhanh cuối tiết.
\end{taisaocoich}
\begin{goiybien}
Có thể đổi thành: một điều bất ngờ, một điều rõ ra, một điều còn muốn hỏi tiếp.
\end{goiybien}
\ngancach

\section{Nếu cho tiết học này một điểm vì độ hữu ích, em cho mấy điểm?}
\begin{cauhoibox}{Câu hỏi}
Nếu cho tiết học này một điểm vì độ hữu ích, em cho mấy điểm? Vì sao?
\end{cauhoibox}
\begin{taisaocoich}
Học sinh thích cho điểm. Câu này vừa vui, vừa cho thầy cô phản hồi thật, lại dễ mở thêm một câu hỏi ngắn phía sau.
\end{taisaocoich}
\begin{goiybien}
Nếu lớp ngại nói thẳng, cho các em giơ tay theo mức điểm rồi mời vài bạn giải thích.
\end{goiybien}
\ngancach

\section{Tuần này lớp mình nên giữ lại điều gì và bỏ bớt điều gì?}
\begin{cauhoibox}{Câu hỏi}
Tuần này lớp mình nên giữ lại điều gì và bỏ bớt điều gì?
\end{cauhoibox}
\begin{taisaocoich}
Đây là câu hỏi rất hợp giờ sinh hoạt vì nó tạo cân bằng: không chỉ soi lỗi mà còn giữ lại điều tốt.
\end{taisaocoich}
\begin{goiybien}
Viết hai cột trên bảng: \textit{Giữ lại} và \textit{Bỏ bớt}. Mỗi bên lấy 3 ý là đã đủ để lớp có cảm giác mình đang cùng nhau chỉnh nhịp.
\end{goiybien}
\ngancach

\section{Câu nào của hôm nay đáng mang sang ngày mai?}
\begin{cauhoibox}{Câu hỏi}
Câu nào của hôm nay đáng mang sang ngày mai?
\end{cauhoibox}
\begin{taisaocoich}
Thay vì chốt chung chung, câu hỏi này buộc học sinh chọn ra một điều đủ đáng nhớ để giữ lại. Nhờ đó cuối tiết bớt rơi.
\end{taisaocoich}
\begin{goiybien}
Có thể cho lớp nói thành cụm ngắn hoặc viết lên giấy nhỏ dán ở góc bảng trước khi ra về.
\end{goiybien}
\ngancach

\section{Tuần sau mình cần nhắc nhau điều gì đầu tiên?}
\begin{cauhoibox}{Câu hỏi}
Tuần sau mình cần nhắc nhau điều gì đầu tiên khi gặp lại?
\end{cauhoibox}
\begin{taisaocoich}
Câu hỏi này làm chiếc cầu giữa hai buổi học. Nó giúp giờ chốt tuần không kết thúc rời rạc mà có hướng đi tiếp.
\end{taisaocoich}
\begin{goiybien}
Hợp cuối tuần hoặc cuối chuỗi bài. Chỉ nên lấy một điều thật ngắn để cả lớp dễ nhớ.
\end{goiybien}
